\chapter{Experimental Methods and Laser Diagnostics}\label{ch:diagnostics}

Models must be validated, and real systems operate under extreme conditions that only experiments can probe: combustion is governed by tightly coupled physics and chemistry, so measurements provide the ground truth, the physical insight and the design guidance that no model supplies on its own~\cite{lecture5,lecture6}. This chapter surveys how flame quantities are measured --- from intrusive probes to laser diagnostics --- and how the canonical laminar flame speed is obtained. The scale gap is itself instructive: a laboratory Bunsen burner runs at about \SI{1}{kW}, a Rolls-Royce Trent at roughly \SI{100}{MW}, so laboratory experiments and high-fidelity simulations must bridge five orders of magnitude in power to inform engine design~\cite{lecture5}. The material follows Lectures~5--7, the experimental chapter of Warnatz \emph{et al.}, and the course laboratory exercise~\cite{lecture5,lecture6,lecture7,warnatz2006combustion,AE4262_LabExercise2026}.

\section{What to measure, and the requirements}

A combustion diagnostic may target the boundary conditions, the chemistry (major, minor, particulate and pollutant species, and surface chemistry), the scalar field (temperature, mixture fraction, density, pressure, speed of sound) or the flow field (mean velocity and r.m.s.\ fluctuations, strain rate, vorticity, integral scales and spectra)~\cite{lecture6}. The ideal measurement is accurate, precise, sensitive, quantitative, non-intrusive, suitable for high pressure, and resolved in time, space and wavelength~\cite{lecture6}. Two resolution requirements are fundamental: the probe volume and the sampling interval must both be \emph{smaller} than the length and time scales of the phenomenon of interest, or the measurement averages over the very structure it is meant to capture~\cite{lecture6}. No single technique meets all of these, so the practical art is matching the technique to the question --- ``there is no best diagnostic''~\cite{lecture6}.


Such measurements are usually made on \emph{canonical} flames whose well-defined, highly repeatable geometries isolate a single piece of physics and supply the controlled boundary conditions needed to validate models and build databases: the laminar premixed \textbf{Bunsen} burner, the rich premixed flat-flame \textbf{McKenna} burner, and the laminar-diffusion \textbf{Wolfhard--Parker} slot burner (two parallel fuel/oxidiser streams forming a steady two-dimensional diffusion flame)~\cite{lecture6,lecture5,warnatz2006combustion}. A recurring tool in optical thermometry is the \textbf{Boltzmann population distribution}: in thermal equilibrium the fraction of molecules in an internal state of energy $E_i$ scales as $\exp(-E_i/k_BT)$, so the ratio of signals from two states --- two absorption or LIF lines, or the envelope of a rotational band in Rayleigh or CARS spectra --- depends only on temperature and yields a non-intrusive measurement of it~\cite{lecture6,lecture7}.
\section{Intrusive probes}

The simplest diagnostics are physical probes~\cite{lecture6}. A \textbf{thermocouple} exploits the Seebeck effect --- two dissimilar metals joined at two junctions held at different temperatures generate a voltage proportional to the temperature difference --- but it suffers from thermal-inertia time lag, intrusiveness (it perturbs the flow it measures), and radiative heat loss that biases the reading at flame temperatures. A \textbf{pressure transducer} measures the deflection of a diaphragm via Hooke's law. An \textbf{exhaust-gas analyser} samples the gas through a cooled probe whose walls quench (``freeze'') further reaction, so the frozen sample can be analysed downstream for its composition~\cite{lecture6}. Probes are cheap and direct but intrusive and slow, which motivates the optical methods.

\section{Measuring the laminar flame speed}

Because $\sLz$ is the master parameter of premixed combustion (Chapter~\ref{ch:premixed}), several standard methods exist, split into stationary and propagating types~\cite{lecture5}. In the \textbf{stationary} methods the flame is held fixed in space:
\begin{itemize}[leftmargin=1.5em,itemsep=2pt]
\item \textbf{Bunsen burner.} The conical premixed flame obeys the kinematic balance of Eq.~\eqref{eq:bunsen}: the burning velocity equals the component of the approach flow normal to the front, $\sLz=v_u\sin\alpha$, with $\alpha$ the cone half-angle. Equivalently the area method gives $\sLz=v_u A_B/A_f$ (burner area over flame area). One measures the supplied volumetric flow and mixture composition to get $v_u$, and the cone angle from a flame image, to extract $\sLz$~\cite{lecture5,transcript_l7}.
\item \textbf{Flat-flame (McKenna) burner.} A one-dimensional, burner-stabilised flame with a uniform exit velocity profile gives the burning velocity directly as $U=Q/A$; it is the standard, highly repeatable calibration flame for combustion research~\cite{lecture5}.

\end{itemize}
In the \textbf{propagating} methods the front moves relative to fixed conditions: the \textbf{tube method} tracks a flame propagating through a premixture-filled tube (correcting for its parabolic shape due to buoyancy and friction), and the \textbf{combustion bomb} ignites the mixture at the centre of a rigid sphere and infers $\sLz$ from the flame radius and pressure histories, accounting for the adiabatic compression of the unburnt gas~\cite{lecture5}.

A practical subtlety from the laboratory exercise illustrates the role of $\sLz$: to keep the Bunsen flow laminar ($\Rey\lesssim2300$) while producing a measurable, stable cone (a design target of about $10^\circ$), the burner diameter must be matched to the fuel. Methane, with $\sLz\approx\SI{0.38}{m\,s^{-1}}$, works with a \SIrange{10}{20}{mm} burner, but hydrogen, nearly an order of magnitude faster, would require a turbulent $\Rey\sim10^4$ at the same diameter, so its burner must be shrunk to about \SIrange{3}{3.5}{mm} to stay laminar~\cite{transcript_l7}. Real Bunsen flames also deviate from a perfect cone at the rim (heat loss/quenching, flame detachment) and at the tip (curvature and stretch effects), which must be accounted for in the analysis~\cite{lecture5,transcript_l7}.

\section{Autoignition and flashback diagnostics}

Lecture~6 focuses on the failure modes of premixed combustors, building on the kinetics of Chapter~\ref{ch:kinetics}~\cite{lecture6}. \textbf{Auto-ignition} is spontaneous ignition due to the thermodynamic state alone, with the ignition delay $\tau_{ig}$ the time a mixture needs to react at given $p$ and $T$; \textbf{auto-ignition flashback} occurs when the ignition delay is shorter than the local residence time, $\tau_{ig}<\tau_{res}$, and even a small pocket that lingers long enough can ignite the whole premixer~\cite{lecture6}. This matters most at high pressure: at \SI{1}{atm} the pressure waves that follow ignition cause oscillations that damp out (ramping the equivalence ratio re-stabilises the flame), but at \SI{20}{atm} the effect is stronger and does not damp~\cite{lecture6}. Two further mechanisms specific to swirl-stabilised combustors are introduced here and revisited for hydrogen in Chapter~\ref{ch:hydrogen}: \textbf{combustion-induced vortex breakdown} (CIVB), in which the flame's adverse pressure field decelerates the swirling core flow and moves the stagnation point upstream, allowing flashback even when the turbulent flame speed is well below the axial velocity; and \textbf{thermoacoustic} flashback, in which large pressure/heat-release oscillations drive transient flow reversal~\cite{lecture6}. Together with the core-flow and boundary-layer mechanisms of Chapter~\ref{ch:premixed}, these complete the flashback taxonomy~\cite{lecture5,lecture6}.

\section{Spectroscopy and light--matter interaction}

Optical diagnostics rest on spectroscopy: the measurement of the electromagnetic spectrum produced when matter interacts with or emits radiation~\cite{lecture6}. A spectrometer disperses light into its constituent wavelengths with a diffraction grating and records intensity versus wavelength, from which species concentrations, equivalence ratio and temperature can be inferred~\cite{lecture6}. The underlying physics is quantised: a photon of energy $\Delta E=h\nu$ is \emph{absorbed} when its energy matches the gap between two molecular energy levels, promoting the molecule to an excited state, and a photon is \emph{emitted} when the molecule relaxes~\cite{lecture6,lecture7}. Three processes are distinguished: \emph{stimulated absorption}, \emph{spontaneous emission} (a photon of random direction and phase), and \emph{stimulated emission} (an incident photon triggers a second, identical photon --- the basis of the laser)~\cite{lecture7}. The family tree of diagnostics divides into intrusive probes and optical methods, the latter into non-laser and laser-based, and into spectroscopic and non-spectroscopic techniques~\cite{lecture6}.

\section{Chemiluminescence}

The cheapest optical method exploits the flame's own light. \textbf{Chemiluminescence} is the natural emission of radiation by electronically excited intermediate species --- principally \ch{OH^*}, \ch{CH^*} and \ch{C_2^*} --- created by specific chemical reactions in the flame~\cite{lecture6}. The intensity and spectral shape of these bands carry quantitative information on the local fuel/air ratio, the heat-release rate and the reaction zone, and imaging through an \ch{OH^*} filter at \SI{310}{nm} marks where reaction occurs~\cite{lecture6}. It is cheap and non-intrusive but \emph{line-of-sight integrated}, so it lacks spatial resolution along the optical path~\cite{lecture6}. The lecture demonstration contrasts two emission mechanisms~\cite{transcript_l7}: chemiluminescence produces discrete electronic-transition \emph{bands} (broadened into bell shapes by the finite instrument resolution and the rotational/vibrational structure), whereas \emph{soot} radiates as a hot blackbody, a broad continuum --- which is why a sooty rich flame glows yellow with a smooth spectrum while a clean lean flame shows discrete radical bands. Richer flames emit more \ch{C_2^*} (green Swan bands) and \ch{CH^*}; leaner flames emit more \ch{OH^*}, but at \SI{310}{nm} in the ultraviolet, invisible to the eye~\cite{transcript_l7,vera2014c2swan}. This is also why hydrogen flames are nearly invisible (and a real safety hazard): they emit essentially nothing in the visible, only \ch{OH} in the UV, so detecting them requires a UV-sensitive instrument~\cite{transcript_l7,schefer2009visible_hydrogen}.

\section{Lasers and how they work}

A laser --- Light Amplification by Stimulated Emission of Radiation --- produces light that is \emph{monochromatic} (one wavelength), \emph{coherent} (all photons in phase) and \emph{directional} (a tight, intense beam), in contrast to the broadband, incoherent, divergent light of a lamp~\cite{lecture6,lecture7}. Lasing requires three ingredients~\cite{lecture7}: a \emph{gain medium} (solid, semiconductor, liquid dye or gas) in which stimulated emission occurs; a \emph{pump} that supplies energy to excite the medium; and an optical \emph{cavity} (two mirrors, one partially reflecting) that feeds photons back through the medium to amplify them. The essential condition is \textbf{population inversion} --- more atoms in the upper than the lower level --- so that stimulated emission outpaces absorption and the light is amplified; this is only possible if the upper level is \emph{metastable}, with a lifetime ($10^{-6}$--$10^{-3}$~s) far longer than an ordinary excited state ($\sim10^{-9}$~s), allowing atoms to accumulate there~\cite{lecture7}. Lasers are either continuous-wave or pulsed (pulses shorter than \SI{0.25}{s} with much higher peak power, able to ``freeze'' the fastest events), and the high-power devices used in combustion are Class~4 and hazardous to eyes and skin~\cite{lecture6,lecture7}.

\section{Laser-based techniques}

When a laser beam meets a molecule, photons are absorbed, emitted or scattered, and the different responses underpin different techniques, broadly split into resonant (absorption, fluorescence) and non-resonant (scattering)~\cite{lecture7}.

\paragraph{Absorption (TDLAS).} A tunable diode laser is set to a wavelength absorbed by the target species, and the transmitted intensity follows the Beer--Lambert law
\begin{equation}
\frac{I}{I_0}=\exp(-k_\nu L)=\exp(-n\,\sigma_\nu\,L),
\label{eq:beer-lambert}
\end{equation}
where $\sigma_\nu$ is the absorption cross-section, $n$ the number density of the absorber and $L$ the path length; the attenuation yields concentration and, from the temperature-dependence of $\sigma_\nu$, temperature~\cite{lecture7,hitran}. It is sensitive but line-of-sight integrated.

\paragraph{Laser-induced fluorescence (LIF).} A tunable laser is tuned to an allowed electronic--vibrational--rotational transition of a chosen species (often \ch{OH}); the excited molecules relax and \emph{fluoresce}, and the fluorescence is collected at \SI{90}{\degree} to the beam~\cite{lecture7,crosley1983}. The signal is proportional to the concentration of the absorbing species (and to the quantum yield, the ratio of fluorescence to incident photons, only $\sim0.1\%$),
\begin{equation}
I_{LIF}=C\,\frac{p}{k_B T}\,I_{laser}\,\frac{g_2}{g_1+g_2}\,A_{21}\,f(T)\,X_{abs},
\label{eq:lif}
\end{equation}
with $A_{21}$ the spontaneous-emission coefficient and $f(T)$ the Boltzmann population fraction~\cite{lecture7}. Spreading the beam into a sheet gives \textbf{planar LIF}, a two-dimensional, \emph{spatially resolved} image of the species distribution --- its great advantage over line-of-sight methods for complex turbulent flames~\cite{lecture7,transcript_l7}. As the transcript stresses, the dye laser is tuned by adjusting the cavity to match the \ch{OH} absorption line, the fluorescence is Stokes-shifted to a longer wavelength, and a bandpass filter blocks the laser wavelength (which would otherwise reflect off surfaces and swamp the image) while passing the fluorescence~\cite{transcript_l7}. LIF is resonant, normally pulsed, linear in laser intensity, and essentially \emph{background-free} (no fluorescence, no signal), which makes it far more sensitive than absorption, though quantitative use generally needs calibration~\cite{lecture7}.

\paragraph{Scattering.} Non-resonant scattering depends on the ratio of particle size $d$ to wavelength $\lambda$~\cite{lecture7}. \emph{Rayleigh} scattering ($d\ll\lambda$, elastic) gives temperature and density; \emph{Raman} scattering (inelastic) gives minor-species concentrations; and \emph{Mie} scattering ($d\sim\lambda$, elastic) is used for droplets and sprays. Related laser methods include particle image velocimetry (PIV) and laser Doppler velocimetry (LDV) for velocity, laser-induced incandescence (LII) for soot, and coherent anti-Stokes Raman (CARS) for temperature~\cite{lecture6,lecture7}.

\paragraph{Reconstructing line-of-sight data: the Abel transform.} For an \emph{axisymmetric} flame, a line-of-sight-integrated image (such as chemiluminescence) can be inverted to a two-dimensional radial profile via the \textbf{Abel transform}, allowing a chemiluminescence image to be compared directly with a planar-LIF slice~\cite{transcript_l7}. The technique is valid only under the axisymmetry assumption, which fails for complex turbulent flames where sheet imaging (LIF) is required instead~\cite{transcript_l7}.

\begin{keybox}
\keyhead{Choosing a diagnostic}
Probes are cheap but intrusive and slow; chemiluminescence is cheap and non-intrusive but line-of-sight; absorption is sensitive but path-integrated; LIF is spatially resolved, sensitive and background-free but complex and usually needs calibration. There is no single best technique --- match it to the quantity, the resolution and the conditions~\cite{lecture6,lecture7}.
\end{keybox}
